\section{L-symbols of a group action}
\label{sec:threecocycle_lsymbols}

Let a group $G$ act on a set $\mathsf X$. An $L$-symbol assigns a nonzero
complex number $L^x_{g,h}$ to each point $x$ and each pair $g,h\in G$; it is
compatible with a scalar three-cochain $\omega$ when the three-cocycle equation
of $\omega$ is balanced by the $L$-symbols. For a one-point set compatibility
says that $\omega$ is the coboundary of $L$, so $L$-symbols measure how far a
three-cocycle is from being trivial along the orbits of the action. This
section develops their gauge freedom and normalization, the consequences of
compatibility at fixed points and at involutions, their behavior under complex
conjugation, and the criterion for trivializing an $L$-symbol by a gauge,
following~\cite{GarreRubio2023MPOSym,FrancoRubio2025MPUGauging}.

\begin{definition}[L-symbols and compatibility]
    \label{def:mpug_l_symbols}
    \lean{LSymbol,
        LSymbol.IsCompatible}
    \leanok
    \uses{def:mpug_scalar_cochains}
    Let $G$ act on a set $\mathsf X$. An $L$-symbol is a function
    $L\colon \mathsf X\times G^2\to\C^\times$, written $L^x_{g,h}$.
    It is compatible with a scalar $3$-cochain $\omega$ when, for every
    $x\in\mathsf X$ and $g,h,k\in G$,
    \begin{align}
        L^x_{g,hk}L^x_{h,k}
         & =\omega(g,h,k)L^{k x}_{g,h}L^x_{gh,k}.
        \label{eq:mpug_l_compatibility}
    \end{align}
    The $L$-symbols were introduced by
    \cite{GarreRubio2023MPOSym}; this is the scalar compatibility relation of
    \cite{FrancoRubio2025MPUGauging}. No transitivity or finiteness assumption
    is imposed on the action.
    % Source anchor: arXiv:2502.20257, eq:omega_and_Ls.
\end{definition}

\begin{definition}[Joint scalar gauge of L-symbols]
    \label{def:mpug_l_gauge}
    \lean{ActionTensorGauge,
        LSymbol.gauge}
    \leanok
    \uses{def:mpug_l_symbols, def:mpug_fusion_gauge}
    A gauge one-cochain is a function
    $\gamma\colon G\times\mathsf X\to\C^\times$, written $\gamma^x_g$.
    Together with a scalar $2$-cochain $\beta\colon G^2\to\C^\times$,
    it changes an $L$-symbol by
    \begin{align}
        (\beta,\gamma)\mathbin{\triangleright}L^x_{g,h}
         & =\frac{\gamma^x_{gh}\beta_{g,h}}
            {\gamma^{h x}_g\gamma^x_h}L^x_{g,h}.
        \label{eq:mpug_l_gauge}
    \end{align}
    This is the joint gauge freedom of \cite{FrancoRubio2025MPUGauging},
    written with the point of $\mathsf X$ as a superscript.
    % Source anchor: arXiv:2502.20257, eq:Lsymbgauge (which spells the point
    %   of X as a second subscript); the superscript spelling gamma^x_g is
    %   the paper's own at lines 5908 and 5988--5990.
\end{definition}

\begin{lemma}[Identity and composition of joint scalar gauges]
    \label{thm:mpug_l_gauge_composition}
    \lean{LSymbol.gauge_one,
        LSymbol.gauge_comp}
    \leanok
    \uses{def:mpug_l_gauge}
    The constant gauges leave every $L$-symbol unchanged.
    Successive joint gauges compose by pointwise multiplication:
    \begin{align}
        (\delta,\varepsilon)\mathbin{\triangleright}
        \bigl((\beta,\gamma)\mathbin{\triangleright}L\bigr)
         & = (\beta\delta,\gamma\varepsilon)
        \mathbin{\triangleright}L.
        \label{eq:mpug_l_gauge_composition}
    \end{align}
\end{lemma}
\begin{proof}\leanok
    \uses{def:mpug_l_gauge}
    Substituting (\ref{eq:mpug_l_gauge}) twice and cancelling the intermediate
    factors yields the composition law, and the identity case
    follows by taking all four cochains equal to $1$.
\end{proof}

\begin{definition}[Normalized L-symbols and gauge one-cochains]
    \label{def:mpug_l_normalized}
    \lean{LSymbol.IsNormalized,
        ActionTensorGauge.IsNormalized}
    \leanok
    \uses{def:mpug_l_symbols, def:mpug_l_gauge}
    As in \cite{FrancoRubio2025MPUGauging}, the $L$-symbol is normalized when,
    for every $x\in\mathsf X$ and $g\in G$,
    \begin{align}
        L^x_{g,e} & =1.
        \label{eq:mpug_l_normalized_right}
    \end{align}
    \begin{align}
        L^x_{e,g} & =1.
        \label{eq:mpug_l_normalized_left}
    \end{align}
    Alongside (\ref{eq:mpug_l_normalized_right}) and
    (\ref{eq:mpug_l_normalized_left}), the gauge one-cochain is normalized
    when
    \begin{align}
        \gamma^x_e & =1
        \label{eq:mpug_action_gauge_normalized}
    \end{align}
    for every $x\in\mathsf X$. This last condition is not a displayed
    equation of the source: it is the scalar form of its standing
    convention that all data involving the identity are trivial.
    % Source anchor: arXiv:2502.20257, eq:triv_Ls for the L-symbol axes;
    %   gamma^x_e = 1 is not a labelled source equation but the standing
    %   convention at lines 1936--1937 that all data involving e are
    %   trivial, applied to the gauge scalar of eq:scalar_act_ten.
\end{definition}

\begin{lemma}[Joint scalar gauges preserve L-symbol compatibility]
    \label{thm:mpug_l_gauge_preservation}
    \lean{LSymbol.IsCompatible.gauge,
        LSymbol.gauge_apply_right_one,
        LSymbol.gauge_apply_left_one,
        LSymbol.isNormalized_gauge_iff,
        LSymbol.IsNormalized.gauge}
    \leanok
    \uses{def:mpug_fusion_gauge, def:mpug_l_symbols, def:mpug_l_gauge, def:mpug_l_normalized}
    If $L$ is compatible with $\omega$, then
    $(\beta,\gamma)\mathbin{\triangleright}L$ is compatible with
    $\beta\mathbin{\triangleright}\omega$. On the identity axes, the gauge
    formula is
    \begin{align}
        ((\beta,\gamma)\mathbin{\triangleright}L)^x_{g,e}
         & =\frac{\beta_{g,e}}{\gamma^x_e}L^x_{g,e}.
        \label{eq:mpug_l_gauge_right}
    \end{align}
    \begin{align}
        ((\beta,\gamma)\mathbin{\triangleright}L)^x_{e,g}
         & =\frac{\beta_{e,g}}{\gamma^{g x}_e}L^x_{e,g}.
        \label{eq:mpug_l_gauge_left}
    \end{align}
    Thus the gauged $L$-symbol is normalized exactly when
    $\beta_{g,e}L^x_{g,e}=\gamma^x_e$ and
    $\beta_{e,g}L^x_{e,g}=\gamma^{g x}_e$ for all $x$ and $g$. In particular,
    normalized $\beta$, $\gamma$, and $L$ give a normalized gauged $L$-symbol.
\end{lemma}
\begin{proof}\leanok
    \uses{def:mpug_fusion_gauge, def:mpug_l_symbols, def:mpug_l_gauge, def:mpug_l_normalized}
    With $L'=(\beta,\gamma)\mathbin{\triangleright}L$, substituting
    (\ref{eq:mpug_l_gauge}) into the left side of
    (\ref{eq:mpug_l_compatibility}) for $L'$ cancels the factor
    $\gamma^x_{hk}$ and leaves
    \begin{align}
        {L'}^x_{g,hk}{L'}^x_{h,k}
         & =\frac{\gamma^x_{ghk}}
            {\gamma^{(hk)x}_g\gamma^{k x}_h\gamma^x_k}\,
        \beta_{g,hk}\beta_{h,k}\,L^x_{g,hk}L^x_{h,k}.
        \label{eq:mpug_l_gauge_left_side}
    \end{align}
    The same substitution on the right side cancels
    $\gamma^{k x}_{gh}$ and leaves
    \begin{align}
        {L'}^{k x}_{g,h}{L'}^x_{gh,k}
         & =\frac{\gamma^x_{ghk}}
            {\gamma^{h(k x)}_g\gamma^{k x}_h\gamma^x_k}\,
        \beta_{g,h}\beta_{gh,k}\,L^{k x}_{g,h}L^x_{gh,k}.
        \label{eq:mpug_l_gauge_right_side}
    \end{align}
    The action law $h(k x)=(hk)x$ makes the two prefactors in $\gamma$ of
    (\ref{eq:mpug_l_gauge_left_side}) and
    (\ref{eq:mpug_l_gauge_right_side}) equal. Applying
    (\ref{eq:mpug_l_compatibility}) for $L$ and $\omega$ to the right of
    (\ref{eq:mpug_l_gauge_left_side}) therefore gives
    \begin{align}
        {L'}^x_{g,hk}{L'}^x_{h,k}
         & =\frac{\beta_{g,hk}\beta_{h,k}}
            {\beta_{g,h}\beta_{gh,k}}\,\omega(g,h,k)\,
                                                     {L'}^{k x}_{g,h}{L'}^x_{gh,k},
        \label{eq:mpug_l_gauge_compatibility}
    \end{align}
    whose scalar factor is
    $(\beta\mathbin{\triangleright}\omega)(g,h,k)$ by
    (\ref{eq:mpug_coboundary}) and
    (\ref{eq:mpug_gauge_action}). Setting either group argument equal
    to $e$ in (\ref{eq:mpug_l_gauge}) gives
    (\ref{eq:mpug_l_gauge_right}) and
    (\ref{eq:mpug_l_gauge_left}), from which the normalization statements
    follow.
\end{proof}

\begin{lemma}[Normalization and involution signs]
    \label{thm:asymex_l_symbol_normalization}
    \lean{LSymbol.IsCompatible.isNormalized_gauge,
        LSymbol.IsCompatible.apply_involution}
    \leanok
    \uses{def:mpug_l_symbols, def:mpug_l_gauge, def:mpug_l_normalized}
    Let $L$ be compatible with a normalized three-cochain $\omega$ for an
    action of $G$ on $X$. The gauge one-cochain
    $\gamma_{g,x}=L^x_{e,e}$, with trivial two-cochain, makes $L$
    normalized. As in the normalization and sign relations used in
    \cite[Section~6]{GarreRubio2023MPOSym}, for every three-cochain $\omega$,
    every normalized $L$ compatible with $\omega$, every $x\in X$, and every
    $g\in G$ with $g^2=e$, one has
    $L^x_{g,g}=\omega(g,g,g)L^{g\cdot x}_{g,g}$.
\end{lemma}
\begin{proof}
    \leanok
    \uses{thm:mpug_l_gauge_preservation, thm:mpug_block_factorization, thm:mpug_relative_block_character}
    Compatibility at $(g,e,e)$ and $(e,e,g)$ gives
    $L^x_{g,e}=L^x_{e,e}/\omega(g,e,e)$ and
    $L^x_{e,g}=\omega(e,e,g)L^{g\cdot x}_{e,e}$. Since $\omega$ is normalized,
    these are $L^x_{g,e}=L^x_{e,e}$ and $L^x_{e,g}=L^{g\cdot x}_{e,e}$.
    The gauge formula then makes both identity-axis values equal to one.
    Compatibility at $(g,g,g)$ reduces to the stated sign relation.
\end{proof}

\begin{lemma}[L-symbols of a family of matrices]
    \label{thm:mpug_l_symbols_action_matrices}
    \lean{LSymbol.isCompatible_one_of_mul_eq_smul,
        LSymbol.isNormalized_of_mul_eq_smul}
    \leanok
    \uses{def:mpug_l_symbols, def:mpug_l_normalized}
    Let $G$ act on $\mathsf X$, and let $S_{g,x}$ be nonzero square matrices
    and $L^x_{g,h}$ nonzero scalars with
    \begin{align}
        S_{g,hx}S_{h,x} & =L^x_{g,h}S_{gh,x}
        \label{eq:mpug_action_matrix_l_symbols}
    \end{align}
    for all $g,h\in G$ and $x\in\mathsf X$, as in
    \cite[lines~4363--4368]{FrancoRubio2025MPUGauging}. Then $L$ is compatible
    with the trivial three-cochain. If in addition $S_{e,x}=\Id$ for every $x$,
    then $L$ is normalized.
\end{lemma}
\begin{proof}\leanok
    \uses{def:mpug_l_symbols}
    Both bracketings of $S_{g,hkx}S_{h,kx}S_{k,x}$ are evaluated with
    (\ref{eq:mpug_action_matrix_l_symbols}). One gives
    $L^{kx}_{g,h}L^x_{gh,k}S_{ghk,x}$ and the other gives
    $L^x_{g,hk}L^x_{h,k}S_{ghk,x}$. Since $S_{ghk,x}\neq0$, the scalars agree.
    With $S_{e,x}=\Id$, the relation at $(g,e)$ and at $(e,g)$ reads
    $S_{g,x}=L^x_{g,e}S_{g,x}$ and $S_{g,x}=L^x_{e,g}S_{g,x}$.
\end{proof}

\begin{lemma}[A point fixed by a subgroup trivializes the restricted three-cocycle]
    \label{thm:mpug_fixed_block_trivializes}
    \lean{LSymbol.isTrivialGaugeClass_comap_of_isCompatible,
        LSymbol.isTrivialGaugeClass_of_isCompatible_of_forall_smul_eq}
    \leanok
    \uses{def:mpug_l_symbols, def:mpug_cochain_restriction, def:mpug_fusion_gauge, def:mpug_trivial_gauge_class}
    Let $L$ be compatible with $\omega$, let $f\colon H\to G$ be a group
    homomorphism, and let $x\in\mathsf X$ satisfy $f(a)x=x$ for every
    $a\in H$. Then $f^*\omega$ has trivial gauge class; explicitly,
    $f^*\omega=d\lambda$ with $\lambda(a,b)=L^x_{f(a),f(b)}$. In particular,
    if some $x$ is fixed by all of $G$, then $\omega$ has trivial gauge class.
    This provides the scalar step both for the statement of
    \cite[lines~740--745]{GarreRubio2023MPOSym} that a subgroup fixing a point
    trivializes the three-cocycle and for the statement of
    \cite[line~1925]{FrancoRubio2025MPUGauging} that a single fixed point
    reduces (\ref{eq:mpug_l_compatibility}) to the three-coboundary condition.
\end{lemma}
\begin{proof}\leanok
    \uses{def:mpug_l_symbols, def:mpug_fusion_gauge, def:mpug_cochain_restriction, def:mpug_trivial_gauge_class}
    For $a,b,c\in H$, with $g=f(a)$, $h=f(b)$, and $k=f(c)$, the invariance
    $kx=x$ turns (\ref{eq:mpug_l_compatibility}) into
    \begin{align}
        \omega(g,h,k)
         & =\frac{L^x_{g,hk}L^x_{h,k}}{L^x_{g,h}L^x_{gh,k}},
    \end{align}
    which is $(d\lambda)(a,b,c)$ by (\ref{eq:mpug_coboundary}) and the
    multiplicativity of $f$. Thus
    $f^*\omega=\lambda\mathbin{\triangleright}1$. For a point fixed by all
    of $G$, taking $f$ to be the identity completes the proof.
\end{proof}

\begin{lemma}[The anomaly at an involution]
    \label{thm:mpug_involution_anomaly_ratio}
    \lean{LSymbol.eq_div_of_isCompatible_of_mul_self_eq_one}
    \leanok
    \uses{def:mpug_l_symbols}
    If $L$ is compatible with $\omega$ and $g^2=e$, then for every
    $x\in\mathsf X$ the formula of
    \cite[line~5121]{FrancoRubio2025MPUGauging} holds:
    \begin{align}
        \omega(g,g,g)
         & =\frac{L^x_{g,g}L^x_{g,e}}{L^{gx}_{g,g}L^x_{e,g}}.
    \end{align}
\end{lemma}
\begin{proof}\leanok
    Setting $h=k=g$ in (\ref{eq:mpug_l_compatibility}) and using $gg=e$
    gives the ratio.
\end{proof}

\begin{lemma}[The involution ratio from the self-pair and identity L-symbols]
    \label{thm:mpug_involution_anomaly_identity_ratio}
    \lean{LSymbol.div_eq_mul_of_isCompatible_of_mul_self_eq_one}
    \leanok
    \uses{def:mpug_l_symbols}
    If $L$ is compatible with $\omega$ and $g^2=e$, then for every
    $x\in\mathsf X$
    \begin{align}
        \frac{L^x_{g,g}L^x_{e,e}}{L^{gx}_{g,g}L^{gx}_{e,e}}
         & =\omega(g,g,g)\,\omega(g,e,g).
        \notag
    \end{align}
    For L-symbols and cocycles that are trivial when an argument is the
    identity, this recovers the relation $L_A/L_B=\omega$ of
    \cite[Section~III.C]{GarreRubioSchuch2024DomainWall} in the form
    $L^x_{g,g}/L^{gx}_{g,g}=\omega(g,g,g)$.
\end{lemma}
\begin{proof}\leanok
    \uses{thm:mpug_involution_anomaly_ratio}
    Compatibility at $(g,e,g)$ gives $L^x_{e,g}=\omega(g,e,g)L^{gx}_{g,e}$,
    and at $(g,e,e)$ it gives $L^z_{e,e}=\omega(g,e,e)L^z_{g,e}$ for $z=x$ and
    $z=gx$. Substituting into the formula of
    Lemma~\ref{thm:mpug_involution_anomaly_ratio} gives the claim.
\end{proof}

\begin{lemma}[Normalized $L$-symbol ratio]
    \label{thm:mpug_sigma_l_ratio}
    \lean{LSymbol.sigma_eq_div}
    \leanok
    \uses{def:mpug_l_symbols, def:mpug_l_normalized, def:mpug_inversion_cochains}
    Let $G$ act on a set $\mathsf X$. Suppose that a normalized $L$-symbol
    $L$ is compatible with a scalar $3$-cochain $\omega$. Then, as in
    \cite[lines~5327--5329]{FrancoRubio2025MPUGauging}, for every
    $x\in\mathsf X$ and $g\in G$,
    \begin{align}
        \sigma(g)
         & =\frac{L^x_{g^{-1},g}}{L^{g x}_{g,g^{-1}}}.
        \label{eq:mpug_sigma_l_ratio}
    \end{align}
    % Source anchor: arXiv:2502.20257, lines 5327--5329.
\end{lemma}
\begin{proof}\leanok
    \uses{def:mpug_l_symbols, def:mpug_l_normalized, def:mpug_inversion_cochains}
    Applying the compatibility relation
    (\ref{eq:mpug_l_compatibility}) to $(g,g^{-1},g)$, the two factors with an
    identity argument are $1$ by normalization, leaving
    (\ref{eq:mpug_sigma_l_ratio}).
\end{proof}

\begin{lemma}[The interchange product of L-symbols]
    \label{thm:mpug_domain_wall_interchange}
    \lean{LSymbol.prod_div_eq_inv_cyclicInvariant,
        LSymbol.prod_Ico_div_eq_inv_cyclicInvariant}
    \leanok
    \uses{def:mpug_l_symbols, def:mpug_cyclic_invariant}
    Let $L$ be compatible with a scalar $3$-cochain $\omega$. For every
    $g\in G$ with $g^n=e$ and every $x\in\mathsf X$,
    \begin{align}
        \prod_{k=0}^{n-1}\frac{L^{gx}_{g,g^k}}{L^x_{g,g^k}}
         & =\Bigl(\prod_{k=0}^{n-1}\omega(g,g^k,g)\Bigr)^{-1}.
        \notag
    \end{align}
    The factor $k=0$ is $L^{gx}_{g,e}/L^x_{g,e}$, which is one for normalized
    L-symbols, so for those the product over $1\leq k<n$ has the same value.
    This is the scalar step in the interchange phase of domain walls
    of~\cite[Section~IV.C]{GarreRubioSchuch2024DomainWall}.
\end{lemma}
\begin{proof}\leanok
    \uses{def:mpug_l_symbols}
    Compatibility (\ref{eq:mpug_l_compatibility}) at $(g,g^k,g)$ gives
    \begin{align}
        \frac{L^{gx}_{g,g^k}}{L^x_{g,g^k}}
         & =\omega(g,g^k,g)^{-1}
        \frac{L^x_{g,g^{k+1}}}{L^x_{g,g^k}}
        \Bigl(\frac{L^x_{g^{k+1},g}}{L^x_{g^k,g}}\Bigr)^{-1},
        \notag
    \end{align}
    and both quotients telescope to one over $k<n$ because $g^n=e$.
\end{proof}

\begin{lemma}[The L-symbol ratio on the stabilizer of two points]
    \label{thm:mpug_domain_wall_projective}
    \lean{LSymbol.ratioCocycle,
        LSymbol.ratioCocycle_isMulCocycle₂}
    \leanok
    \uses{def:mpug_l_symbols, def:is_cocycle}
    Let $L$ be compatible with a scalar $3$-cochain $\omega$, let
    $y,z\in\mathsf X$, and let $H_{y,z}$ be the subgroup of elements fixing
    both $y$ and $z$. Then
    \begin{align}
        \alpha(a,b)
         & =\frac{L^y_{a,b}}{L^z_{a,b}}
        \notag
    \end{align}
    is a two-cocycle on $H_{y,z}$. In~\cite[Section~IV.C]{GarreRubioSchuch2024DomainWall}
    this is the factor system of the projective action on domain walls of the
    subgroup fixing every point, which is contained in $H_{y,z}$.
\end{lemma}
\begin{proof}\leanok
    \uses{def:mpug_l_symbols}
    For $a,b,c\in H_{y,z}$ we have $cy=y$ and $cz=z$, so dividing the
    compatibility relation (\ref{eq:mpug_l_compatibility}) at $y$ by the one at
    $z$ cancels $\omega(a,b,c)$ and gives the two-cocycle equation for
    $\alpha$.
\end{proof}

\begin{lemma}[Time reversal makes the anomaly and the L-symbols jointly real]
    \label{thm:mpug_time_reversal_real_l_symbols}
    \lean{LSymbol.star_gauge,
        LSymbol.gauge_modulusUnit_pos,
        LSymbol.star_gauge_phaseRoot,
        LSymbol.exists_gauge_star_eq_self}
    \leanok
    \uses{def:mpug_fusion_gauge, def:mpug_l_gauge, def:mpug_phase_root}
    Let $\omega$ be a scalar $3$-cochain, $L$ an $L$-symbol, $\beta$ a
    scalar $2$-cochain and $\gamma$ a gauge one-cochain with
    $\beta\mathbin{\triangleright}\omega=\overline\omega$ and
    $(\beta,\gamma)\mathbin{\triangleright}L=\overline L$. Let $s$ and $t$ be
    the unimodular square roots of the phases of $\beta$ and $\gamma$. Then
    $s\mathbin{\triangleright}\omega$ and
    $(s,t)\mathbin{\triangleright}L$ are both real-valued. This result is not stated
    in~\cite[Section~4.3]{GarreRubio2023MPOSym}.
    The source derives $\beta\mathbin{\triangleright}\omega=\overline\omega$
    and, with the two-cochain omitted,
    $(1,\gamma)\mathbin{\triangleright}L=\overline L$; the hypothesis here,
    with the same $\beta$ in both relations, is the gauge-consistent pair
    and is assumed rather than derived from time reversal.
    % Source anchor: arXiv:2203.12563, lines 873--891 and 935--957.
\end{lemma}
\begin{proof}\leanok
    \uses{def:mpug_l_gauge, thm:mpug_l_gauge_composition, thm:mpug_time_reversal_real_anomaly, thm:mpug_phase_root}
    The claim for $\omega$ is Lemma~\ref{thm:mpug_time_reversal_real_anomaly}.
    For $L$, writing $\beta=s^2N$ and $\gamma=t^2M$ with
    $N_{g,h}=|\beta_{g,h}|$ and $M^x_g=|\gamma^x_g|$ gives, as in that proof,
    $\overline{(s,t)\mathbin{\triangleright}L}
        =(N,M)\mathbin{\triangleright}((s,t)\mathbin{\triangleright}L)$, whose
    prefactor
    $M^x_{gh}N_{g,h}/(M^{h x}_gM^x_h)$ is a positive real number, hence
    equal to one by Lemma~\ref{thm:mpug_phase_root}.
\end{proof}

\begin{lemma}[The L-symbol relation imposed by time reversal]
    \label{thm:mpug_time_reversal_l_symbol_relation}
    \lean{LSymbol.star_apply_eq_gauge_one_apply,
        LSymbol.phase_sq_eq_of_gauge_eq_star}
    \leanok
    \uses{def:mpug_l_gauge}
    Let $T\neq0$ and $P$ be vectors in a complex vector space and let
    $L^x_{g,h}$ and $\gamma$ be as in Definition~\ref{def:mpug_l_gauge}.
    If $P=\overline{L^x_{g,h}}\,T$ and
    $\gamma^x_h\gamma^{h x}_gP=L^x_{g,h}\gamma^x_{gh}T$, then
    \begin{align}
        \frac{\overline{L^x_{g,h}}}{L^x_{g,h}}
         & =\frac{\gamma^x_{gh}}{\gamma^x_h\gamma^{h x}_g}.
        \label{eq:mpug_time_reversal_l_symbol}
    \end{align}
    More generally, if $(\beta,\gamma)\mathbin{\triangleright}L=\overline L$,
    then for all $x$, $g$ and $h$,
    \begin{align}
        \left(\frac{L^x_{g,h}}{|L^x_{g,h}|}\right)^2
         & =\frac{\gamma^{h x}_g\gamma^x_h}{\gamma^x_{gh}\beta_{g,h}}.
    \end{align}
    In~\cite[Section~4.3]{GarreRubio2023MPOSym}, the first hypothesis is
    the complex conjugate of the relation defining $L$, and the second is
    the relation derived from time-reversal symmetry. The source prints the consequence as
    $|L^x_{g,h}|^2=\gamma^x_{gh}/(\gamma^x_h\gamma^{h x}_g)$; that form fails
    for an on-site $\mathbb Z_2\times\mathbb Z_2$ symmetry, and
    (\ref{eq:mpug_time_reversal_l_symbol}) is the local correction. The second display is the
    corrected sense in which the $L$-symbols squared are gauge
    transformations: the square of the phase of $L$ is, while its modulus is
    unconstrained.
    % Source anchor: arXiv:2203.12563, lines 935--957 (TRSF1group at
    %   line 954).
    % Local fix: https://github.com/LionSR/TNLean/blob/main/docs/paper-gaps/glm23_time_reversal_l_symbol.tex
\end{lemma}
\begin{proof}\leanok
    \uses{def:mpug_l_gauge}
    Substituting the first hypothesis into the second gives
    $\gamma^x_h\gamma^{h x}_g\overline{L^x_{g,h}}\,T
        =L^x_{g,h}\gamma^x_{gh}T$, and $T\neq0$ gives
    (\ref{eq:mpug_time_reversal_l_symbol}). For the second claim,
    $\overline L=cL$ with
    $c=\gamma^x_{gh}\beta_{g,h}/(\gamma^{h x}_g\gamma^x_h)$, so
    $|L|^2=\overline LL=cL^2$ and $(L/|L|)^2=c^{-1}$.
\end{proof}

\begin{definition}[Block scalar and block independence]
    \label{def:mpug_block_independence_factor}
    \lean{LSymbol.ell,
        LSymbol.IsBlockIndependentEll}
    \leanok
    \uses{def:mpug_l_symbols}
    Following~\cite{FrancoRubio2025MPUGauging}, the points of $\mathsf X$
    are called blocks. For an $L$-symbol, define
    \begin{align}
        \ell^x_{a,b;r}
         & =\frac{L^x_{ar,r^{-1}b}}{L^x_{a,b}}.
        \label{eq:mpug_block_scalar}
    \end{align}
    The scalar $\ell$ is block-independent when
    \begin{align}
        \ell^x_{a,b;r}
         & =\ell^y_{a,b;r}
        \label{eq:mpug_block_independence}
    \end{align}
    for all $a,b,r\in G$ and $x,y\in\mathsf X$.
    % Source anchor: arXiv:2502.20257, app:block_indep, lines 6965--6974.
\end{definition}

\begin{definition}[Four scalar gauge formulations]
    \label{def:mpug_block_independence_gauge_forms}
    \lean{LSymbol.IsBlockConstant,
        LSymbol.IsTrivialEll,
        LSymbol.IsTrivial,
        LSymbol.HasTrivialGauge,
        LSymbol.HasBlockConstantGauge,
        LSymbol.HasTrivialEllGauge,
        LSymbol.IsBlockIndependent}
    \leanok
    \uses{def:mpug_l_gauge, def:mpug_block_independence_factor}
    Writing $L'=(\beta,\gamma)\mathbin{\triangleright}L$ for a joint scalar
    gauge of an $L$-symbol, the four scalar $L$-symbol properties used in the
    proof of the source proposition are, respectively, the existence of $\beta$
    and $\gamma$ such that
    \begin{align}
        (L')^x_{g,h}       & =1
                           &                     & \text{for all $x,g,h$},
        \label{eq:mpug_block_form_trivial}                                     \\
        (L')^x_{a,b}       & =(L')^y_{a,b}
                           &                     & \text{for all $x,y,a,b$},
        \label{eq:mpug_block_form_constant}                                    \\
        \ell[L']^x_{a,b;r} & =1
                           &                     & \text{for all $x,a,b,r$},
        \label{eq:mpug_block_form_trivial_ell}                                 \\
        \ell[L']^x_{a,b;r} & =\ell[L']^y_{a,b;r}
                           &                     & \text{for all $x,y,a,b,r$}.
        \label{eq:mpug_block_form_independent_ell}
    \end{align}
    % Source anchor: arXiv:2502.20257, prop:technical01, lines 6975--6989.
\end{definition}

\begin{lemma}[Corrected block factorization]
    \label{thm:mpug_block_factorization}
    \lean{LSymbol.IsCompatible.apply_right_one,
        LSymbol.blockFactor,
        LSymbol.factorization_of_isBlockIndependentEll}
    \leanok
    \uses{def:mpug_l_symbols, def:mpug_block_independence_factor}
    Suppose that $L$ is compatible with $\omega$ and that $\ell$ is
    block-independent. Fix $x_0\in\mathsf X$ and set
    \begin{align}
        A(g,h)
         & =\bigl(\omega(gh,e,e)\ell^{x_0}_{g,h;h}\bigr)^{-1}.
        \label{eq:mpug_block_factor}
    \end{align}
    Then, with the reciprocal factor recorded in
    \cite{gap:fbc25_block_independence_reciprocal_errors}, for every
    $x\in\mathsf X$ and $g,h\in G$,
    \begin{align}
        L^x_{g,h}
         & =A(g,h)L^x_{e,e}.
        \label{eq:mpug_corrected_block_factorization}
    \end{align}
    % Local correction:
    % docs/paper-gaps/fbc25_block_independence_reciprocal_errors.tex.
    % Source anchor: arXiv:2502.20257, eq:Lfact, lines 6994--7002.
\end{lemma}
\begin{proof}\leanok
    \uses{def:mpug_block_independence_factor}
    Definition~\ref{def:mpug_block_independence_factor} gives
    \begin{align}
        L^x_{g,h}
         & =\frac{L^x_{gh,e}}{\ell^x_{g,h;h}}.
        \label{eq:mpug_block_factorization_first}
    \end{align}
    Block independence replaces $x$ by $x_0$ in the denominator.
    Compatibility at $(gh,e,e)$ gives
    $L^x_{gh,e}=L^x_{e,e}/\omega(gh,e,e)$, which proves
    (\ref{eq:mpug_corrected_block_factorization}).
\end{proof}

\begin{lemma}[Relative block character]
    \label{thm:mpug_relative_block_character}
    \lean{LSymbol.IsCompatible.apply_left_one,
        LSymbol.relativeCharacter,
        LSymbol.relativeCharacter_mul_base,
        LSymbol.relativeCharacter_mul}
    \leanok
    \uses{thm:mpug_block_factorization}
    Under the hypotheses of Lemma~\ref{thm:mpug_block_factorization}, define
    \begin{align}
        \rho(g)
         & =\bigl(\ell^{x_0}_{e,g;g}\omega(g,e,e)
        \omega(e,e,g)\bigr)^{-1}.
        \label{eq:mpug_relative_block_character}
    \end{align}
    As recorded in \cite{gap:fbc25_block_independence_reciprocal_errors}, the
    reciprocal factor $(\ell^{x_0}_{e,g;g})^{-1}$ corrects the corresponding
    printed factor. Then
    \begin{align}
        L^{g x}_{e,e}
         & =\rho(g)L^x_{e,e},
        \label{eq:mpug_relative_block_transport} \\
        \rho(gh)
         & =\rho(g)\rho(h).
        \label{eq:mpug_relative_block_multiplicative}
    \end{align}
    Thus $\rho\colon G\to\C^\times$ is a one-dimensional character. No
    transitivity or finiteness assumption on the action is required.
    % Local correction:
    % docs/paper-gaps/fbc25_block_independence_reciprocal_errors.tex.
    % Source anchor: arXiv:2502.20257, eq:defrho, lines 7003--7018.
\end{lemma}
\begin{proof}\leanok
    \uses{thm:mpug_block_factorization}
    Compatibility at $(e,e,g)$ gives
    $L^x_{e,g}=\omega(e,e,g)L^{g x}_{e,e}$, and comparing this with
    (\ref{eq:mpug_corrected_block_factorization}) at $(e,g)$ yields
    (\ref{eq:mpug_relative_block_transport}). Applying that identity first
    to $h$ and then to $g$, and comparing with its application to $gh$, proves
    (\ref{eq:mpug_relative_block_multiplicative}) by cancellation.
\end{proof}

\begin{lemma}[Explicit inverse gauge]
    \label{thm:mpug_block_independence_inverse_gauge}
    \lean{LSymbol.factorFusionGauge,
        LSymbol.factorActionGauge,
        LSymbol.gauge_one_eq_of_isBlockIndependentEll,
        LSymbol.gauge_inv_eq_one_of_isBlockIndependentEll,
        LSymbol.hasTrivialGauge_of_isBlockIndependentEll}
    \leanok
    \uses{def:mpug_l_gauge, thm:mpug_l_gauge_composition, thm:mpug_block_factorization, thm:mpug_relative_block_character}
    With $A$ and $\rho$ as above, define the forward gauges by
    \begin{align}
        \beta^+_{g,h}
         & =\frac{A(g,h)}{\rho(h)},
        \label{eq:mpug_forward_fusion_gauge} \\
        \gamma^{+,x}_g
         & =\bigl(L^{g x}_{e,e}\bigr)^{-1}.
        \label{eq:mpug_forward_action_gauge}
    \end{align}
    Their joint gauge multiplier equals $L$:
    \begin{align}
        (\beta^+,\gamma^+)\mathbin{\triangleright}1
         & =L.
        \label{eq:mpug_forward_gauge_multiplier}
    \end{align}
    Consequently the pointwise inverse gauges trivialize $L$:
    \begin{align}
        ((\beta^+)^{-1},(\gamma^+)^{-1})
        \mathbin{\triangleright}L
         & =1.
        \label{eq:mpug_inverse_gauge_trivialization}
    \end{align}
    As recorded in \cite{gap:fbc25_block_independence_reciprocal_errors}, the
    distinction between the forward multiplier and its pointwise inverse
    corrects the orientation of the displayed source calculation.
    If $\mathsf X$ is nonempty, a choice of $x_0$ therefore gives a
    trivializing joint scalar gauge for every compatible block-independent
    $L$-symbol.
    % Local correction:
    % docs/paper-gaps/fbc25_block_independence_reciprocal_errors.tex.
    % Source anchor: arXiv:2502.20257, eq:Lgauge and lines 7019--7027.
\end{lemma}
\begin{proof}\leanok
    \uses{thm:mpug_l_gauge_composition, thm:mpug_block_factorization, thm:mpug_relative_block_character}
    Substituting (\ref{eq:mpug_forward_fusion_gauge}) and
    (\ref{eq:mpug_forward_action_gauge}) into
    (\ref{eq:mpug_l_gauge}), equation
    (\ref{eq:mpug_relative_block_transport}) reduces the multiplier to
    $A(g,h)L^x_{e,e}$, which equals $L^x_{g,h}$ by
    (\ref{eq:mpug_corrected_block_factorization}). Joint gauges compose by
    pointwise multiplication, so the inverse gauges give
    (\ref{eq:mpug_inverse_gauge_trivialization}).
\end{proof}

\begin{lemma}[Trivial and block-constant gauges]
    \label{thm:mpug_trivial_iff_block_constant}
    \lean{LSymbol.IsTrivial.isBlockConstant,
        LSymbol.IsBlockConstant.isBlockIndependentEll,
        LSymbol.HasTrivialGauge.hasBlockConstantGauge,
        LSymbol.HasBlockConstantGauge.isBlockIndependent,
        LSymbol.hasTrivialGauge_iff_hasBlockConstantGauge}
    \leanok
    \uses{def:mpug_block_independence_gauge_forms, thm:mpug_l_gauge_composition, thm:mpug_l_gauge_preservation, thm:mpug_block_independence_inverse_gauge}
    Let $L$ be an $L$-symbol compatible with a scalar $3$-cochain,
    and suppose that $\mathsf X$ is nonempty. A joint scalar gauge trivializes
    $L$ if and only if a joint scalar gauge makes $L^x_{g,h}$ independent of
    $x$.
    % Source anchor: arXiv:2502.20257, prop:technical01, items 1 and 2.
\end{lemma}
\begin{proof}\leanok
    \uses{thm:mpug_l_gauge_composition, thm:mpug_l_gauge_preservation, thm:mpug_block_independence_inverse_gauge}
    A trivial $L$-symbol is block-constant. Conversely, a block-constant
    representative has block-independent $\ell$, so applying
    Lemma~\ref{thm:mpug_block_independence_inverse_gauge} to that
    representative and composing the two joint gauges trivializes $L$.
\end{proof}

\begin{lemma}[Trivial gauges and trivial block scalars]
    \label{thm:mpug_trivial_iff_trivial_ell}
    \lean{LSymbol.IsTrivial.isTrivialEll,
        LSymbol.IsTrivialEll.isBlockIndependentEll,
        LSymbol.HasTrivialGauge.hasTrivialEllGauge,
        LSymbol.HasTrivialEllGauge.isBlockIndependent,
        LSymbol.hasTrivialGauge_iff_hasTrivialEllGauge}
    \leanok
    \uses{def:mpug_block_independence_gauge_forms, thm:mpug_l_gauge_composition, thm:mpug_l_gauge_preservation, thm:mpug_block_independence_inverse_gauge}
    Let $L$ be an $L$-symbol compatible with a scalar $3$-cochain,
    and suppose that $\mathsf X$ is nonempty. A joint scalar gauge trivializes
    $L$ if and only if a joint scalar gauge makes every
    $\ell^x_{a,b;r}$ equal to $1$.
    % Source anchor: arXiv:2502.20257, prop:technical01, items 1 and 3.
\end{lemma}
\begin{proof}\leanok
    \uses{thm:mpug_l_gauge_composition, thm:mpug_l_gauge_preservation, thm:mpug_block_independence_inverse_gauge}
    A trivial $L$-symbol has $\ell=1$. Conversely, the identity $\ell=1$
    implies block independence, so applying
    Lemma~\ref{thm:mpug_block_independence_inverse_gauge} and composing the
    gauges trivializes $L$.
\end{proof}

\begin{lemma}[Scalar gauge triviality and block independence]
    \label{thm:mpug_trivial_iff_block_independent}
    \lean{LSymbol.IsBlockIndependent.hasTrivialGauge,
        LSymbol.hasTrivialGauge_iff_isBlockIndependent}
    \leanok
    \uses{def:mpug_block_independence_gauge_forms, thm:mpug_l_gauge_composition, thm:mpug_l_gauge_preservation, thm:mpug_block_independence_inverse_gauge}
    Let $L$ be an $L$-symbol compatible with a scalar $3$-cochain,
    and suppose that $\mathsf X$ is nonempty. A joint scalar gauge trivializes
    $L$ if and only if a joint scalar gauge makes $\ell^x_{a,b;r}$ independent
    of $x$.
    % Source anchor: arXiv:2502.20257, prop:technical01, items 1 and 4.
\end{lemma}
\begin{proof}\leanok
    \uses{thm:mpug_l_gauge_composition, thm:mpug_l_gauge_preservation, thm:mpug_block_independence_inverse_gauge}
    Gauge triviality implies block constancy and hence block independence.
    For the converse, apply
    Lemma~\ref{thm:mpug_block_independence_inverse_gauge} to the
    block-independent representative and compose its trivializing gauge with
    the original gauge.
\end{proof}

\begin{lemma}[Four scalar gauge formulations of block independence]
    \label{thm:mpug_block_independence_four_forms}
    \lean{LSymbol.hasTrivialGauge_tfae}
    \leanok
    \uses{def:mpug_block_independence_gauge_forms, thm:mpug_trivial_iff_block_constant, thm:mpug_trivial_iff_trivial_ell, thm:mpug_trivial_iff_block_independent}
    Let $L$ be an $L$-symbol compatible with a scalar $3$-cochain,
    and suppose that $\mathsf X$ is nonempty. The four scalar gauge
    properties (\ref{eq:mpug_block_form_trivial}),
    (\ref{eq:mpug_block_form_constant}),
    (\ref{eq:mpug_block_form_trivial_ell}), and
    (\ref{eq:mpug_block_form_independent_ell}) of $L$ are equivalent. This is
    the scalar core of the block-independence proposition of
    \cite{FrancoRubio2025MPUGauging}, whose fourth condition is stated there
    for the data that produce the $L$-symbol.
    % Source anchor: arXiv:2502.20257, prop:technical01, lines 6975--6989.
\end{lemma}
\begin{proof}\leanok
    \uses{thm:mpug_trivial_iff_block_constant, thm:mpug_trivial_iff_trivial_ell, thm:mpug_trivial_iff_block_independent}
    Lemmas~\ref{thm:mpug_trivial_iff_block_constant},
    \ref{thm:mpug_trivial_iff_trivial_ell}, and
    \ref{thm:mpug_trivial_iff_block_independent} identify each of the last
    three properties with the first.
\end{proof}
